\section{Evaluation Design}\label{sec:evaluation}

\subsection{Data}

We use the audio-visual speech portion of the Ryerson Audio-Visual Database of Emotional Speech and Song (RAVDESS; \citealp{livingstone2018}). Twenty-four professional actors (12 female, 12 male) speak two lexically matched, neutral statements in eight emotions: neutral, calm, happy, sad, angry, fearful, disgust and surprised. All emotions except neutral are produced at two intensities. This gives 1440 clips (60 per actor), 96 of them neutral and 192 for each other class. Clips last 3.69\,s on average (SD 0.34\,s) and are recorded at 1280$\times$720 pixels and 29.97 frames per second. The statements carry no emotional content, so the linguistic channel is uninformative by design and the study isolates the paralinguistic and facial channels. RAVDESS is acted, frontal and recorded under studio conditions. This is a limitation for external validity, but it is an advantage for a controlled study of degradations, because the clean condition is truly clean and every degradation can be injected and measured.

For the confirmatory evaluation of Cycle~5 we use a subset of the Crowd-sourced Emotional Multimodal Actors Dataset (CREMA-D; \citealp{cao2014}). Its 91 actors (43 female, 48 male) are more diverse in age and ethnicity than those of RAVDESS and were recorded at 480$\times$360 pixels against a uniform backdrop. We use the clips of four of the twelve sentences (``Don't forget a jacket'', ``I think I have a doctor's appointment'', ``The airplane is almost full'', ``I wonder what this is about'') with unspecified intensity, in all six emotions of the corpus (neutral, happy, sad, angry, fearful, disgust). This gives 2178 clips, 363 per class and 18--24 per actor, lasting 2.53\,s on average (SD 0.48\,s). The six classes are a subset of the RAVDESS classes, which allows a cross-corpus test. All processing steps and model settings are identical to those used for RAVDESS.

\subsection{Feature extraction}

Audio is decoded to 16\,kHz mono. Video is sampled at about 5 frames per second (every sixth frame; 18.9 frames per clip on average). A YuNet face detector \citep{wu2023yunet} locates the largest face in each frame (resized to 640$\times$360). The face box is enlarged by 15\% and cropped from the full-resolution frame. The detector returned a face in every sampled frame in all conditions, so the fallback crops were never needed. The four extractors are:
\begin{itemize}
\item \textbf{A0 prosody} (cheap): 280 descriptors computed with librosa \citep{mcfee2015}. These are means and standard deviations of 40 MFCCs with their first and second derivatives, energy, zero-crossing rate, spectral centroid, bandwidth, roll-off and contrast, chroma means, and fundamental-frequency statistics (YIN, 65--400\,Hz).
\item \textbf{A1 wav2vec~2.0} (expensive): the self-supervised \texttt{wav2vec2-base} model \citep{baevski2020}, pre-trained on unlabelled read English speech with no emotion supervision. Hidden states of all 13 layers are averaged over time and then over layers (768-d).
\item \textbf{V0 MobileNetV3} (cheap): the ImageNet-trained \texttt{mobilenet\_v3\_small} \citep{howard2019}. The 576-d pooled embedding of each face crop is summarised by its mean and standard deviation over frames (1152-d).
\item \textbf{V1 CLIP} (expensive): the image encoder of CLIP ViT-B/32 \citep{radford2021}. L2-normalised 512-d frame embeddings are summarised by their mean and standard deviation (1024-d). We also use CLIP zero-shot (V1-ZS). Each class is represented by the averaged text embedding of three prompts (e.g.\ ``a photo of a \{angry\} face.''), and frame-level similarities are averaged over the clip.
\end{itemize}
None of the pre-trained models was trained with supervised emotion labels. CLIP's web-scale image--text pre-training may include emotion vocabulary and similar images, which is why we also test it zero-shot.

\subsection{Stream degradations}

In Cycles~1--3 models are trained on clean data. In Cycle~4 they are trained on clean data plus one degraded copy, as described in Section~\ref{sec:envelope}. Test data are presented clean or under one of eight degradations that mimic common streaming impairments \citep{hendrycks2019}. On audio we add white Gaussian noise at a signal-to-noise ratio of 20, 10 or 0\,dB (\emph{snr20}, \emph{snr10}, \emph{snr0}). On video we apply three impairments to the full frame before face detection. \emph{compress} downscales to 320$\times$180, encodes with JPEG quality 8 and upscales again, emulating a low-bit-rate stream. \emph{lowlight} applies a gamma of 2.2, scales brightness by 0.35 and adds sensor noise ($\sigma=6$). \emph{blur} applies a Gaussian blur ($21\times21$ kernel, $\sigma=6$). Two combined conditions impair both channels: \emph{snr10+compress} and \emph{snr0+lowlight}. Noise is seeded by clip identifier, so every condition is exactly reproducible. The confirmatory evaluation adds three \emph{realistic} degradations that were not used in Cycles~1--4. \emph{cafe5} and \emph{street5} add recorded noise from a cafeteria and from a street with traffic (DEMAND corpus, channel 1; \citealp{thiemann2013}), from a random segment, at 5\,dB SNR. Unlike white noise, these noises contain babble, clatter and low-frequency rumble. \emph{h264} re-encodes the video with the H.264 codec \citep{wiegand2003} at 320$\times$180 pixels and 100\,kbit/s, which produces the blocking and smearing typical of a congested video call. On CREMA-D we use the audio conditions snr10, snr0, cafe5 and street5, the video conditions compress, lowlight and h264, and the two combined conditions.

\subsection{Protocol, models and statistics}

Evaluation is strictly subject-independent (DR1). We use six outer folds, each holding out four actors (two female, two male). Every clip is predicted exactly once by a model that never saw its speaker. Within each training split, the regularisation strength $C\in\{0.01,0.1,1\}$ of each logistic regression \citep{pedregosa2011} is selected by a four-fold inner cross-validation grouped by actor. The same inner predictions are used to fit the temperatures, the cascade thresholds and the monitor thresholds. No test person is used for any fitting or selection. On CREMA-D we use seven outer folds of 13 actors each, balanced by sex as far as possible, with the same inner procedure. For the cross-corpus test, models are trained on all clips of the six shared classes in one corpus (1056 RAVDESS clips or all 2178 CREMA-D clips) and tested on the other corpus, and they are compared with six-class models evaluated within each corpus.

We report accuracy and macro-averaged F1 (macro-F1) pooled over the six test folds, with 95\% confidence intervals from a cluster bootstrap over actors (2000 resamples). Systems are compared on the 24 per-actor accuracies with the Wilcoxon signed-rank test \citep{wilcoxon1945,demsar2006}, and Holm's correction is applied within families of comparisons. Calibration is measured by the expected calibration error with 15 bins (ECE; \citealp{naeini2015}). Group parity is assessed by comparing accuracy and ECE between female and male actors.

\subsection{Cost measurement}

Each processing step was timed per clip (wall-clock time) on the clean condition: audio decoding, video decoding with face detection, and each extractor. The logistic-regression heads take microseconds and are not included. Timings were taken on an Intel Xeon processor at 2.1\,GHz with two PyTorch threads per worker, without a GPU. Table~\ref{tab:cost} reports the means over 1380 clips (the 60 clips of one actor, processed during a pilot with a different thread setting, are excluded). We express cost in CPU seconds per clip and as a real-time factor, which is compute time divided by media duration. The cost of a cascade or gating policy is computed from these component means and the components each clip actually runs. Two workers ran in parallel on four cores, so the timings include some resource contention.

\subsection{Streaming simulation}

To evaluate monitoring (DR4) we simulate streams for each outer fold. A stream presents the 240 test clips of the fold in random order under clean conditions and then the same clips, in a new random order, under one of the eight degradations. Each stream is repeated with 20 random seeds, giving 120 streams per condition. A recovery scenario presents 160 clean, 160 degraded and 160 clean clips. Control streams of 480 clean clips, the 240 clean clips in two random orders, measure false alarms under random variation. Speaker-blocked control streams present the same clean clips grouped by actor, in two different actor orders. They test whether natural speaker changes trigger alarms. For the end-to-end evaluation (Section~\ref{sec:results}), the cost of every clip is the sum of the mean costs of the components that clip actually runs, plus the monitor computation. We compare two label-free signals. The first is the stage-1 confidence, monitored with a one-sided Page--Hinkley test for a fall in confidence ($\delta=0.005$, $\lambda=5$; \citealp{page1954}) and with ADWIN ($\delta=0.002$). The second is the per-channel input monitor score, monitored with ADWIN ($\delta=0.002$). We report the detection rate, the detection delay in clips after the change point, the false alarms per 1000 clean clips, and whether the first alarm comes from the degraded channel. Cycle~5 compares the monitor variants on three scenarios with the same seeds. S1 is the clean-then-degraded stream described above, now including the realistic degradations. S4 contains speaker-blocked clean streams of 480 clips. On RAVDESS these are the fold's four actors in two actor orders, as above. On CREMA-D they consist of 20 sessions of 24 clips, drawn from the 13 test actors in random order, so that some speakers appear twice. S5 is S4 with a degradation that starts at a random clip between 120 and 360, possibly in the middle of a session, and persists to the end. Here a \emph{false mark} is a change of a channel's state to degraded while that channel is clean, and we report false marks per 1000 clips and the share of streams with at least one false mark. For gating, these scenarios are evaluated with the monitor variant chosen by the pre-specified rule of Section~\ref{sec:qd}. To test whether the monitors notice a change of corpus, we also fit them on one corpus and score clean clips of the other, and we run streams that switch from 240 held-out clips of the training corpus to 240 clips of the other corpus.
